\documentclass[aps,prl,twocolumn,superscriptaddress]{revtex4-2}
\usepackage{amsmath,amssymb,graphicx,xcolor}
\usepackage{amsthm} % for the proposition / observation environments
\newtheorem{observation}{Observation}
\newtheorem{proposition}{Proposition}
% gvq: renders its second argument only; the first argument is a key of
% the project's numeric registry (gate C21).
\newcommand{\gvq}[2]{#2}

\begin{document}

\title{The Fine Structure Constant as a Spectral Coincidence on the Three-Sphere}

\author{J.~Loutey}
\affiliation{Independent Researcher, Kent, Washington}

\date{February 25, 2026}

\begin{abstract}
We record a numerical observation. The number
$K = \pi(42 + \pi^2/6 - 1/40) = 137.036064$ lies within
$4.8\times10^{-7}$ (relative) of the inverse fine structure constant
$\alpha^{-1}$, and the smallest positive root of the cubic
$\alpha^3 - K\alpha + 1 = 0$ lies within $8.8\times10^{-8}$ of $\alpha$
(CODATA 2018). The three ingredients are spectral data of the unit
three-sphere $S^3$: $B = 42$ is the degeneracy-weighted trace of the
$SO(3)$ Casimir $l(l+1)$ over the scalar $S^3$ harmonics of the first
three shells, the cutoff being the only one at which $B/N = \dim S^3 = 3$;
$F = \zeta(2) = \pi^2/6$ is the Dirichlet series of the shell degeneracy
$n^2$ at $s = 4$; and $\Delta = 1/40$ is the reciprocal of a Dirac
degeneracy on $S^3$. The combination rule $K = \pi(B + F - \Delta)$ is an
Observation, not a derivation: its weights, signs, prefactor and the
cubic are chosen, and we list which steps are computed and which are
chosen. In a search over $N = 1.92\times10^9$ formulas built from
ingredients of the same kind, ten match $\alpha$ more closely
($p = 5.2\times10^{-9}$ per formula), so the expected chance count is
$Np \approx 10$: within that search the formula is not unusual. The
search was run after the formula was found and does not correct for
that. Twelve mechanisms proposed to generate all three ingredients from
one spectral construction were tested and eliminated; we tabulate them.
\end{abstract}

\maketitle

%======================================================================
\section{Introduction}
%======================================================================

The fine structure constant $\alpha = e^2/(4\pi\epsilon_0\hbar c)
\approx 1/137.036$ is an unexplained input to the Standard
Model~\cite{feynman1985}. Numerical formulas for it have a long and
unsuccessful history~\cite{eddington1935,barrow2002}; Wyler's
symmetric-space expression~\cite{wyler1969} reached comparable
precision and was criticized for lacking any mechanism connecting the
mathematics to the coupling~\cite{robertson1971}.

This paper records one such formula and states what is and is not
known about it. Its ingredients come from the unit three-sphere $S^3$,
onto which Fock's 1935 stereographic projection maps the hydrogen
bound-state problem~\cite{fock1935,paper7}. The GeoVac label lattice
has as nodes the labels $(n,l,m)$ of the $S^3$ harmonics, on which the
Laplace--Beltrami operator acts as the diagonal $SO(4)$ Casimir
$n^2-1$ (the lattice's edge Laplacian does not converge to that
operator, Paper~7). Nothing in the formula uses the edge structure; it
uses only the labels, their degeneracies, and two spectral series.

We state the status at the outset. The combination rule
$K = \pi(B + F - \Delta)$ is a numerical Observation. Each ingredient
has a closed form, but the way they are combined is chosen, and a
look-elsewhere count shows that a search
of this size is expected to produce about ten matches this close by
chance. An earlier version of this work~\cite{alpha_hopf_v1} used a
``symplectic impedance'' ratio with one fitted parameter
($6.2\times10^{-5}$ relative error); it is superseded.

The sections below define the three ingredients and the cutoff, state
the formula and the cubic, give the look-elsewhere count, separate
computed from chosen steps, record where each ingredient lives and the
twelve eliminated mechanisms, and state the scope of any
Hopf-fibration reading of the formula.


%======================================================================
\section{The three ingredients}\label{sec:ingredients}
%======================================================================

\subsection{$B$: a Casimir trace over the scalar $S^3$ harmonics}

Each shell $n$ of the scalar harmonics on $S^3$ carries the $(j,j)$
representation of $SO(4) \cong (SU(2)\times SU(2))/\mathbb{Z}_2$,
$j = (n-1)/2$, which branches under the rotation subgroup $SO(3)$ into
$l = 0, \ldots, n-1$. Define the degeneracy-weighted trace of the
$SO(3)$ Casimir over the first $n_{\max}$ shells,
\begin{equation}\label{eq:casimir}
  B(n_{\max}) = \sum_{n=1}^{n_{\max}} \sum_{l=0}^{n-1}
  (2l+1)\,l(l+1).
\end{equation}
At $n_{\max} = 3$,
\begin{align}
  B(3) &= \underbrace{0}_{(1,0)}
         + \underbrace{0}_{(2,0)} + \underbrace{6}_{(2,1)}
         + \underbrace{0}_{(3,0)} + \underbrace{6}_{(3,1)}
         + \underbrace{30}_{(3,2)} \notag \\
       &= \gvq{b_casimir}{42},
\end{align}
with subscripts the $(n,l)$ pairs.

$B$ is an $S^3$ quantity, not an $S^2$ one. The Laplacian on $S^2$ has
eigenvalues $l(l+1)$ with multiplicity $2l+1$, and its trace over
$l \le 2$ is $0 + 6 + 30 = 36$; $B$ is larger because the $l = 1$
multiplet occurs in two shells ($n = 2$ and $n = 3$). Its value also
depends on how each shell is labelled (see the scope section below).

The inner sum has the closed form
\begin{equation}\label{eq:innersum}
  \sum_{l=0}^{n-1}(2l+1)\,l(l+1) = \frac{n^2(n^2-1)}{2},
\end{equation}
so the per-shell average of the $SO(3)$ Casimir equals the $SO(4)$
Casimir eigenvalue,
\begin{equation}\label{eq:casimir_identity}
  \langle C_3 \rangle_n
  = \frac{1}{n^2}\sum_{l=0}^{n-1}(2l+1)\,l(l+1)
  = \frac{n^2-1}{2} = C_4(n),
\end{equation}
and
\begin{equation}\label{eq:Bclosed}
  B(m) = \frac{m(m+1)(2m+1)(m+2)(m-1)}{20}.
\end{equation}

\subsection{The cutoff $n_{\max} = 3$}

With $N(m) = \sum_{n=1}^{m} n^2 = m(m+1)(2m+1)/6$ the number of states
through shell $m$,
\begin{equation}\label{eq:ratioclosed}
  \frac{B(m)}{N(m)} = \frac{3(m+2)(m-1)}{10}.
\end{equation}
Setting $B/N = \dim(S^3) = 3$ gives $(m+2)(m-1) = 10$, i.e.\
$(m+4)(m-3) = 0$, whose only positive root is $m = 3$
(Table~\ref{tab:selection}). The criterion $B/N = \dim S^3$ is chosen;
given it, the cutoff is computed. Equivalently, by
Eq.~(\ref{eq:casimir_identity}), the degeneracy-weighted average of the
$SO(4)$ Casimir over shells $1 \le n \le n_{\max}$ equals $\dim S^3$.

\begin{table}[tb]
\centering
\begin{tabular}{cccc}
\hline\hline
$n_{\max}$ & $B(n_{\max})$ & $N(n_{\max})$ & $B/N$ \\
\hline
1 & 0   & 1  & 0.000 \\
2 & 6   & 5  & 1.200 \\
3 & 42  & 14 & 3.000 \\
4 & 162 & 30 & 5.400 \\
5 & 462 & 55 & 8.400 \\
\hline\hline
\end{tabular}
\caption{$B/N = 3$ holds only at $n_{\max} = 3$,
Eq.~(\ref{eq:ratioclosed}).}
\label{tab:selection}
\end{table}

\emph{Verification.} \texttt{tests/test\_paper2\_ingredients.py}
recomputes $B(3) = 42$ term by term from Eq.~(\ref{eq:casimir}), proves
Eqs.~(\ref{eq:innersum})--(\ref{eq:ratioclosed}) symbolically, and
checks that $B/N = 3$ holds at $n_{\max} = 3$ and fails at
$n_{\max} \in \{1,2,4,5\}$, all in exact arithmetic.

\emph{$S^3$ specificity of the criterion.} The analogous ratio built
from eigenvalues alone is not specific to $S^3$. For the round sphere
$S^d$ with $\lambda_k = -k(k+d-1)$ and degeneracies $g_k$, the formal
moment $B_{\mathrm{formal}}(m) = \tfrac12\sum_{k=0}^{m} g_k|\lambda_k|$
satisfies $B_{\mathrm{formal}}(m)/N(m) = d$ exactly at $m = 2$ for every
$d \ge 1$~\cite{phase6_hopf} (\texttt{tests/test\_paper2\_corrections.py}).
The $B$ of Eq.~(\ref{eq:casimir}) is the trace of the $SO(3)$
sub-Casimir over the $SO(4) \to SO(3)$ branching, which equals
$B_{\mathrm{formal}}$ only for $d = 3$, where the diagonal
$SO(3) \subset SO(4)$ has embedding index one.

\subsection{$F$: a Dirichlet series of the shell degeneracy}

The $n$-th scalar shell of $S^3$ has degeneracy $n^2$. Its Dirichlet
series at $s = 4$ is
\begin{equation}\label{eq:F_dirichlet}
  F \;=\; \sum_{n\geq 1}\frac{n^2}{n^{4}}
    \;=\; \zeta_R(2) \;=\; \frac{\pi^2}{6} \;=\; \gvq{f_fock_dirichlet}{1.6449}\ldots
\end{equation}
The exponent $s = 4$ is chosen; it equals the maximum degree $d_{\max}$
of the packing construction (Paper~0) and the dimension of the ambient
$\mathbb{R}^4$. The same number is half the spectral zeta function of
the circle of circumference $2\pi$ at $s = 1$:
$\zeta_{S^1}(s) = 2\sum_{k\ge1}k^{-2s}$, so
$F = \tfrac12\zeta_{S^1}(1)$, the $k > 0$ modes. Either reading is a
choice of which spectral series to evaluate and where.

\subsection{$\Delta$: a reciprocal degeneracy}

Two closed forms give $\Delta = 1/40$ at the cutoff. The first combines
the scalar eigenvalue of the last kept shell with the number of states
below it,
\begin{equation}\label{eq:boundary}
  \Delta = \frac{1}{|\lambda_3|\, N(2)}
         = \frac{1}{8 \times 5} = \gvq{delta_dirac}{1/40},
\end{equation}
with $\lambda_n = -(n^2-1)$. The second is the reciprocal of a Dirac
degeneracy: the Camporesi--Higuchi spectrum on unit $S^3$ has
$|\lambda_k^{\mathrm{D}}| = k + 3/2$ with multiplicity
$g_k^{\mathrm{D}} = 2(k+1)(k+2)$, $k = 0, 1, \ldots$, and
\begin{equation}\label{eq:Delta_dirac}
  \Delta^{-1} \;=\; g_3^{\mathrm{D}} \;=\; 2 \cdot 4 \cdot 5 \;=\; 40 .
\end{equation}
As functions of the cutoff $m$, $|\lambda_m|N(m-1)$ and
$g_m^{\mathrm{D}}$ are different polynomials; among positive integers
they agree only at $m = 3$. The Dirac index here is the
Camporesi--Higuchi one: $k = 3$ is the fourth Dirac level
($|\lambda_3^{\mathrm{D}}| = 9/2$), so the shared integer $3$ is shared
across two labelling conventions. The physical significance of $\Delta$
is not established.


%======================================================================
\section{The formula and the cubic}\label{sec:formula}
%======================================================================

Combining the three ingredients,
\begin{equation}\label{eq:K}
\begin{split}
  K &= \pi\!\left(B + F - \Delta\right)
     = \pi\!\left(42 + \frac{\pi^2}{6} - \frac{1}{40}\right)\\
    &= 137.036064\ldots,
\end{split}
\end{equation}
which differs from $\alpha^{-1}_{\mathrm{CODATA\,2018}} = 137.035999084$
by $4.8\times10^{-7}$ relative. The physical $\alpha$ approximately
satisfies (to $8.8\times10^{-8}$ relative through the cubic's root; an
Observation) the relation $1/\alpha + \alpha^2 = K$, i.e.\ the depressed
cubic
\begin{equation}\label{eq:cubic}
  \alpha^3 - K\alpha + 1 = 0 .
\end{equation}
Its discriminant $4K^3 - 27$ is positive, so it has three real roots;
we take the smallest positive one, the perturbative one, which gives
$\alpha^{-1} = 137.036011$. Both $K$ and the root lie above the
measured value, the root closer by a factor of about five. With CODATA
2022 ($137.035999177$) the two relative differences are
$4.76\times10^{-7}$ and $8.75\times10^{-8}$.


\emph{Verification.} \texttt{tests/test\_paper2\_corrections.py}
recomputes the root ($1/\alpha = 137.036011$, above CODATA, within
$10^{-7}$ relative).

The cubic is not selected by any structure we know. Every depressed
cubic $s^3 - Ks + 1$ is the characteristic polynomial of a traceless
$3\times3$ circulant with $bc = K/3$ and $b^3 + c^3 = -1$
(\texttt{tests/test\_paper2\_ingredients.py}), so that realization fixes
nothing; the self-consistency form $1/\alpha + \alpha^2 = K$ is a
choice.

\begin{proposition}[The cubic's root is transcendental]
\textbf{[SYMBOLIC]} The physical fine structure constant only
\emph{approximately} satisfies the cubic of Eq.~(\ref{eq:cubic}) (to
$8.8\times10^{-8}$ relative, far outside the experimental uncertainty),
so the cubic is not an exact law, and this statement concerns the
cubic's smallest positive root $\alpha_K$, not the physical $\alpha$:
$\alpha_K$ is algebraic over $\mathbb{Q}(\pi)$ and transcendental over
$\mathbb{Q}$. Nothing follows for the physical $\alpha$, and the
combination rule remains an Observation.
\end{proposition}
\emph{Proof.} From $\alpha_K^3 - K\alpha_K + 1 = 0$ and $\alpha_K\neq0$,
$K = \alpha_K^2 + 1/\alpha_K$. If $\alpha_K$ were algebraic, so would
be $K$. But $K = \pi(B+F-\Delta) = \tfrac{1679}{40}\,\pi +
\tfrac{1}{6}\,\pi^{3}$ is a nonconstant polynomial in $\pi$ with
rational coefficients, hence transcendental (Lindemann). $\square$
(The mechanism is specific:\ roots of polynomials with transcendental
coefficients \textit{can} be algebraic --- e.g.\ $x^2 - \pi x + (\pi-1)
= (x-1)(x-\pi+1)$ has the root $x = 1$ --- but that needs the
$\pi$-dependence to cancel at the root.  That cannot happen here:\ in
$x^3 - Kx + 1$ the transcendental $\pi$ enters only through the single
coefficient $K$, the other coefficients being rational, so an algebraic
root $x_0 \neq 0$ would force $K = (x_0^3 + 1)/x_0$ to be algebraic;\ in
the example $\pi$ enters two coefficients, and they cancel at $x = 1$.)
(\texttt{tests/test\_paper2\_cubic\_root\_transcendence.py}.)


%======================================================================
\section{Look-elsewhere count}\label{sec:statistical}
%======================================================================

A formula built from $\pi$, $42$ and $\pi^2/6$ that matches $\alpha^{-1}$
this closely could be a numerical accident. To estimate how often such
matches arise, we ran a combinatorial search~\cite{phase2_audit}.

\subsection{Search space}

The search considers $\alpha^{-1} = T(A + B + C)$ and the cubic
$\alpha^3 - T(A+B+C)\alpha + 1 = 0$, where
\begin{itemize}
\item $A$, $B$, $C$ are drawn from a pool of spectral quantities:
  eigenvalues $|\lambda_n|$, degeneracies $n^2$ and $2l+1$, Casimir
  values $l(l+1)$, cumulative counts $N(n)$, their pairwise products
  up to~50, signed ratios of these, and the transcendentals $\zeta(2)$,
  $\zeta(3)$, $\zeta(4)$, $\ln 2$, $\pi$, $\pi^2$ (743 distinct term
  values);
\item $T$ ranges over 14 prefactors: $\pi$, $2\pi$, $\pi^2$, $\pi/2$,
  $4\pi$, $4\pi^2$, $1/(2\pi)$, $1$, $2$, $1/2$, $\pi/4$, $\pi/6$,
  $1/\pi$, $\sqrt{\pi}$.
\end{itemize}
Both readings (linear and cubic) give $N = 1.92\times10^9$ formulas in
all. The pool was built to contain this formula's own ingredients, so
the count is conditional on that pool.

\subsection{Results}

\begin{table}[tb]
\centering
\begin{tabular}{rrr}
\hline\hline
Threshold & Hits & Fraction \\
\hline
$10^{-2}$          & 1{,}283{,}592  & $6.7 \times 10^{-4}$ \\
$10^{-3}$          &   125{,}027    & $6.5 \times 10^{-5}$ \\
$10^{-4}$          &    14{,}476    & $7.5 \times 10^{-6}$ \\
$10^{-5}$          &       682      & $3.5 \times 10^{-7}$ \\
$10^{-6}$          &        37      & $1.9 \times 10^{-8}$ \\
$10^{-7}$          &        11      & $5.7 \times 10^{-9}$ \\
$8.8\times 10^{-8}$ &       10      & $5.2 \times 10^{-9}$ \\
\hline\hline
\end{tabular}
\caption{Formulas of the search matching $\alpha^{-1}$ within each
relative threshold. The fraction is roughly linear in the threshold
from $10^{-2}$ to $10^{-4}$; below that the counts are small and include
repeated values (one number reached by several formulas).}
\label{tab:hits}
\end{table}

Ten formulas match $\alpha^{-1}$ more closely than the cubic's root
($8.81\times10^{-8}$), so this formula ranks eleventh, and the fraction
at its precision is
\begin{equation}
  p = \frac{10}{1.92\times10^{9}} = 5.2\times10^{-9}.
\end{equation}
The expected number of chance matches at least this close in a search
of this size is therefore
\begin{equation}\label{eq:chance_count}
  N p \;=\; 1.92\times10^{9} \times 5.2\times10^{-9} \;\approx\; 10 ,
\end{equation}
and ten are found. The ten take only three distinct values (the
closest is $4\pi(5/49 + 14/15 + \pi^2)$, at $2.1\times10^{-8}$). At the
precision of $K$ itself ($4.8\times10^{-7}$, the linear reading), 27
entries of the search match more closely, one of them this formula's
own cubic reading~\cite{phase2_audit}.

Within this search, then, a match of this quality is what chance
produces about ten times over; $p$ is a fraction per formula, not a
significance level. The search was designed and run after the formula
was found, so it also does not account for the unregistered search
that produced the formula, nor for the choice of pool. The count does
not support the formula and does not refute it; it shows that precision
alone carries no evidential weight here. (The full sweep is not
regression-gated; the driver is preserved at
\texttt{tests/paper2\_support/phase2\_combinatorial.py}.)


%======================================================================
\section{What is computed and what is chosen}\label{sec:ledger}
%======================================================================

Table~\ref{tab:ledger} separates the steps from spectral data to the
number that are computed from those that are chosen. The formula has no
continuous fitted parameter, but it has several discrete choices, and
each of them enters the look-elsewhere count above as one draw from a
pool of alternatives.

\begin{table*}[tp]
\centering
\begin{tabular}{p{0.36\textwidth}p{0.14\textwidth}p{0.42\textwidth}}
\hline\hline
Step & Status & Note \\
\hline
Weight $(2l+1)\,l(l+1)$ in $B$ & chosen & over the rotation labelling of each shell \\
Criterion $B/N = \dim S^3$ & chosen & \\
Cutoff $n_{\max} = 3$ & computed & from the criterion, Eq.~(\ref{eq:ratioclosed}) \\
Values $B = 42$, $F = \pi^2/6$, $\Delta = 1/40$ & computed & closed forms, Eqs.~(\ref{eq:casimir})--(\ref{eq:Delta_dirac}) \\
Series and point for $F$ ($s = 4$, or $s=1$ on $S^1$) & chosen & \\
Form of $\Delta$ & chosen & two forms agree at $m = 3$ only \\
Additive combination & chosen & \\
Signs $(+,+,-)$ & chosen, then checked & the other seven sign patterns miss by $\ge 1.1\times10^{-3}$ \\
Prefactor $\pi$ & chosen & \\
Cubic $1/\alpha + \alpha^2 = K$ & chosen & the circulant realization selects nothing \\
Root of the cubic & computed & the smallest positive root \\
\hline\hline
\end{tabular}
\caption{Computed and chosen steps of the formula. ``Chosen'' means no
construction we know fixes the step; it is the one that gives the
match.}
\label{tab:ledger}
\end{table*}

The sign check is the one place where the alternatives were scored. Of
the eight patterns $\pi(\pm B \pm F \pm \Delta)$ at $n_{\max} = 3$,
$(+,+,-)$ gives $4.8\times10^{-7}$, the next best $(+,+,+)$ gives
$1.1\times10^{-3}$, and the rest are $7.4\times10^{-2}$ or worse.
Flipping the sign of $\Delta$ moves $K$ by $2\pi\Delta$, which is
$1.1\times10^{-3}$ relative, and $\pi(42 + \pi^2/6)$ alone already lies
within $5.8\times10^{-4}$ of $\alpha^{-1}$; the check therefore amounts
to taking the sign of $\Delta$ that brings the sum closer, and it is not
independent of the match itself.


%======================================================================
\section{Where the ingredients live, and twelve eliminated
mechanisms}\label{sec:mechanisms}
%======================================================================

\begin{proposition}[Three homes]\label{thm:three_homes}
The three ingredients $B$, $F$, $\Delta$ of the combination rule
(an Observation;\ Eq.~\eqref{eq:K}) have the following spectral homes on
$S^3$:
\begin{itemize}
  \item[$(B)$] a finite truncated Casimir trace over the scalar
    Laplace--Beltrami harmonics at the cutoff $m = 3$, with weight
    $(2l+1)\,l(l+1)$, Eq.~(\ref{eq:casimir});
  \item[$(F)$] an infinite Dirichlet series of the scalar shell
    degeneracy at $s = 4$, $F = \sum_{n} n^2 \cdot n^{-4} = \zeta_R(2)$;
  \item[$(\Delta)$] the reciprocal of a single-level Dirac degeneracy,
    $\Delta^{-1} = g_3^{\mathrm{D}} = 40$ (Camporesi--Higuchi index).
\end{itemize}
They lie in two spectral sectors of $S^3$ (scalar and spinor) and are
three different kinds of object (finite trace, infinite series,
single-level multiplicity).
\end{proposition}
\noindent The proof is the closed forms of
Eqs.~(\ref{eq:casimir}), (\ref{eq:F_dirichlet}) and (\ref{eq:Delta_dirac}).
\noindent \textbf{[OBSERVATION]} No single spectral construction on $S^3$ generating all
three has been exhibited:\ twelve single-mechanism derivations have been
eliminated (an empirical record, not a theorem;\ cf.\ the E6 scope
remark of Paper~32).

\begin{observation}[Cross-sector sum]\label{obs:cross_sector}
% [retracted 2026-09-29: k-match-inside-theorem] -- the numerical match of
% the combination rule previously sat inside the theorem environment; it is
% an Observation and is stated here, outside the proposition.
The combination rule $K = \pi(B + F - \Delta)$ is a
\emph{cross-sector sum} of three ingredients with separate algebraic
origins. That the sum comes within $4.8\times 10^{-7}$ of
$\alpha^{-1}/\pi$ for $K$ itself (to $8.8\times 10^{-8}$ through the
cubic's root) is a numerical observation, not a consequence of
Proposition~\ref{thm:three_homes}; no common-generator mechanism for it
is known.
\end{observation}

\subsection{The twelve mechanisms}

Table~\ref{tab:mechanisms} lists the mechanisms tested for a single
spectral construction producing $B$, $F$ and $\Delta$ together (or the
combination $K$). Eight come from the structural program of Phases
4B--4I, one from Sprint~A, and three from Sprint~K-CC~\cite{alpha_program_record,alpha_sprint_a_memo}.
The record of Phases 4B--4I counted the discrete and the continuum
circle-fibre tests as two mechanisms (nine in all); the corpus count of
twelve treats them as one, row~4. Elimination means the construction
was computed and missed; it is an empirical record, not an exhaustion.

\begin{table*}[tp]
\centering
\begin{tabular}{rp{0.30\textwidth}p{0.11\textwidth}p{0.47\textwidth}}
\hline\hline
\# & Mechanism & Where & Result \\
\hline
1 & Spectral comparison of $S^3$ with the untwisted $S^1\times S^2$ & 4B, $\alpha$-A &
 zeta values, regularized determinants, truncated Casimir traces and heat-kernel coefficients of the two manifolds reproduce none of $K$, $K/\pi$, $B+F$, $B+F-\Delta$ better than $10^{-3}$ \\
2 & Traces of higher Casimir powers on $S^3$ & 4B, $\alpha$-B &
 $\sum_n n^2 c_2(n)^k$: only the $k=1$ trace gives a component ($42 = B$); none gives $\pi^2/6$ or $K/\pi$ \\
3 & Spectral invariants of the scalar label graph and its $S^2$ quotient at $n_{\max}=3$ & 4B, $\alpha$-D &
 trace, log-determinant, zeta, Cheeger constant and entropy of either Laplacian hit no target unless $(2l+1)l(l+1)$ is reinserted by hand \\
4 & Circle-fibre traces, discrete (Fock weight) and continuum & 4C $\alpha$-E, 4D $\alpha$-H &
 path-graph fibre zetas are rational under every weighting; continuum fibre traces give $\pi\cdot\mathbb{Q}$ at each asymptotic order (Jacobi inversion), never $\pi^2/6$ \\
5 & Spectral geometry of $S^5$ & 4E, $\alpha$-I &
 volumes give integer powers of $\pi$, heat-kernel coefficients rationals, determinants $\zeta(\mathrm{odd})$ and $\log 2$; $\pi^2$ never appears beside a rational \\
6 & Zeta-value combinations for $\Delta$ & 4G, $\alpha$-K &
 3{,}228 candidates at 50 digits: none within $10^{-25}$ of $1/40$ \\
7 & Standard-Model running (one-loop vacuum polarization) as the origin of $\Delta$ & 4H, SM-A--F &
 one-loop QED running gives $\Delta(1/\alpha) = 1/40$ only at $\mu \approx 575$~MeV, no natural scale; $\sum N_cQ_f^2 = 8 = |\lambda_3|$ over three generations ($8/3$ per generation), a coincidence; $1/40$ is no Chern number, $\eta$-invariant or Chern--Simons level of $SU(5)$, $Spin(10)$ or $E_8$ except by renaming \\
8 & Lift of $B$, $F$, $\Delta$ to the Dirac spectrum on $S^3$ & 4I, D2--D4 &
 obstructions 1--3 below \\
9 & Connes--Chamseddine bulk spectral action on $S^3$ & Sprint A, $\alpha$-SP &
 the Dirac heat trace on round $S^3$ has exactly two power terms, $\tfrac{\sqrt{\pi}}{2}t^{-3/2} - \tfrac{\sqrt{\pi}}{4}t^{-1/2}$, with signs $(+,-)$ (Paper~51), so it cannot supply three ingredients with the pattern $(+,+,-)$ \\
10 & A cutoff $\Lambda$ with $\mathrm{Tr}\,e^{-D^2/\Lambda^2} = K/\pi$, identified by integer relation & K-CC (a) &
 PSLQ of $\Lambda = 3.7102454\ldots$ at 100 digits against a 244-element basis: no relation \\
11 & A natural $\Lambda$ from $S^3$ spectral data & K-CC (b) &
 only tautological hits: every target has one real root $\Lambda$ \\
12 & $B$, $F$, $\Delta$ as three terms of one heat-kernel expansion & K-CC (c) &
 the Seeley--DeWitt coefficients of $D^2$ on unit $S^3$ lie in $\sqrt{\pi}\,\mathbb{Q}$, $B$ is a finite trace of a different operator and $\Delta^{-1}$ a single-level count, so no one expansion contains all three \\
\hline\hline
\end{tabular}
\caption{The twelve eliminated mechanisms. ``Where'' names the phase or
sprint and its track. Rows 1--8 are from the record of Phases
4B--4I, row 9 from Sprint~A (2026-04-18), rows 10--12 from Sprint~K-CC
(2026-05-03)~\cite{alpha_program_record,alpha_sprint_a_memo}.}
\label{tab:mechanisms}
\end{table*}

\subsection{The three Dirac obstructions (row 8)}\label{sec:obstructions}

\paragraph{Obstruction 1 (to lifting $B$).}
The scalar Laplacian on $S^3$ has a zero mode at $(n=1, l=0)$, which
forces the factor $(m-1)$ in Eq.~(\ref{eq:Bclosed}). The Dirac operator
has no zero mode ($|\lambda_k^{\mathrm{D}}| \ge 3/2$), and its
cumulative traces,
\begin{align}
  \sum_{k=0}^{M} g_k^{\mathrm{D}}
  &= \frac{2(M+1)(M+2)(M+3)}{3},\\
  \sum_{k=0}^{M} g_k^{\mathrm{Weyl}}
  &= \frac{(M+1)(M+2)(M+3)}{3},
\end{align}
carry no such factor. The single hit
$|\lambda_{2}^{\mathrm{D}}|\cdot g_{2}^{\mathrm{Weyl}} = (7/2)\cdot12 = 42$
follows from $(m-1)(m+2) = 10$ at $m = 3$, an elementary identity.

\paragraph{Obstruction 2 (to lifting $F$).}
The Dirac degeneracy splits as $g_m = 2m^2 + 2m$ in the integer label
$m = k+1$, so its Dirichlet series at $s = 4$ is
\begin{equation}\label{eq:D_dirac_zeta3}
  \sum_{m\geq 1}\frac{2m(m+1)}{m^4}
  \;=\; 2\,\zeta_R(2) + 2\,\zeta_R(3).
\end{equation}
\emph{Conditional on} $\zeta(3)/\zeta(2) \notin \mathbb{Q}$ (open:
Ap\'ery's theorem gives only $\zeta(3) \notin \mathbb{Q}$), no rational
combination of these Dirac-sector series is a rational multiple of $F$;
isolating $F$ means projecting onto the $m^2$ part, i.e.\ returning to
the scalar sector. In the Camporesi--Higuchi index the series is
$\sum_{k\ge0} 2(k+1)(k+2)/(k+3/2)^4 = \pi^2 - \pi^4/12$, which is
unconditionally not a rational multiple of $F$ (Lindemann).

\paragraph{Obstruction 3 (to unifying $\Delta$ with $B$ or $F$).}
Under the $S^1$ action of the Hopf fibration (the Orduz
$S^1$-equivariant decomposition of the spinor bundle), the 40 states of
the Dirac level $k = 3$ split as $20 + 20$, by chirality or,
equivalently, into half-integer and integer charge sectors. The
decomposition produces neither $B$ nor $F$, and the
scalar factor $8 = |\lambda_3|$ of Eq.~(\ref{eq:boundary}) is a scalar
eigenvalue, not a spinor quantity, so the two factorizations of $40$
do not combine.

The base spectrum and degeneracies are in \texttt{geovac/dirac\_s3.py}
(\texttt{tests/test\_dirac\_s3.py}, including $g_3^{\mathrm{D}} = 40$);
the obstruction identities themselves were computed in archived sprint
drivers and are recorded in \texttt{docs/dirac\_s3\_verdict.md}.
The ingredients are also not Seeley--DeWitt coefficients of the
perturbative spectral action (row 12); an earlier ``Seeley--DeWitt
cancellation theorem'' supporting a stronger reading is withdrawn
(2026-10-02; the exact $S^3$ heat traces give coefficient ratios $2$,
$-1$, $0, \ldots$, Paper~18).
% [retracted 2026-10-02: sd-ratio-four]


%======================================================================
\section{Scope of a Hopf-fibration reading}\label{sec:scope}
%======================================================================

Earlier versions of this paper read the three ingredients as data of
the three spaces of the Hopf fibration $S^1 \to S^3 \to S^2$
\cite{hopf1931}: $B$ on the base, $F$ on the fibre, $\Delta$ on the total
space. Two measured facts limit any such reading.

\emph{The labelling matters.} $B$ is computed over the rotation
grouping of each shell (orbits of $SO(3)$, labelled by $l$). Grouping
each shell instead by the chiral $SU(2)$ factor whose torus is the
Hopf fibre gives a Casimir trace of exactly $B/2 = 21$ at the same
cutoff, because $\mathrm{Tr}\,L^2 = \mathrm{Tr}\,A^2$ on each hydrogen
shell. The degeneracy $n^2$, and hence $F$, does not depend on the
grouping (\texttt{tests/test\_paper25\_hopf\_grouping.py}; record
\texttt{docs/predictions/2026-10-31\_hopf\_grouping.md}). So $B = 42$ is
attached to the rotation labelling, not to the fibration.

\emph{The fibration is a one-centre structure.} The Hopf fibre is
generated by $L_z + A_z$. With a second nucleus on the $z$ axis, the
commutant of its potential inside $\mathfrak{so}(4)$ is spanned by
$L_z$ alone, and the fibre generator is broken at relative size $0.38$
(\texttt{tests/test\_atlas\_so4\_commutant.py}; record
\texttt{docs/predictions/}\linebreak[0]\texttt{2026-10-30\_so4\_commutant\_door11.md}). Any
reading of $K$ through the fibration therefore concerns a single
Coulomb centre.


%======================================================================
\section{Comparison}\label{sec:comparison}
%======================================================================

Wyler's formula~\cite{wyler1969} gives $\alpha^{-1} = 137.03608$, about
$6\times10^{-7}$ from the measured value, comparable to $K$ here. Both
formulas rest on a combination rule without a mechanism. What this
paper adds is bookkeeping: the computed-versus-chosen ledger of
Table~\ref{tab:ledger}, a look-elsewhere count (Table~\ref{tab:hits}),
and a record of the mechanisms that were tested and failed
(Table~\ref{tab:mechanisms}).


%======================================================================
\section{Conclusion}
%======================================================================

The number $K = \pi(42 + \pi^2/6 - 1/40)$ lies within
$4.8\times10^{-7}$ of $\alpha^{-1}$, and the smallest positive root of
$\alpha^3 - K\alpha + 1 = 0$ within $8.8\times10^{-8}$ of $\alpha$. The
ingredients are spectral data of $S^3$: a Casimir trace over the scalar
harmonics of the first three shells (the only cutoff with
$B/N = \dim S^3$), a Dirichlet series of the shell degeneracy, and a
reciprocal Dirac degeneracy. The combination rule
$K = \pi(B + F - \Delta)$ is a numerical Observation. Its weights,
signs, prefactor and the cubic are chosen (Table~\ref{tab:ledger}). In a
search of $1.92\times10^9$ formulas of the same kind the expected chance
count at this precision is $Np \approx 10$ ($p = 5.2\times10^{-9}$ per
formula), and ten closer matches are found, so precision alone does not
distinguish this formula. Twelve proposed single-construction mechanisms
were tested and eliminated (Table~\ref{tab:mechanisms}). No derivation
of the combination rule exists.


\begin{acknowledgments}
This work supersedes an earlier ``symplectic impedance''
approach~\cite{alpha_hopf_v1} by the same author. The computational audit
that checked the formula's arithmetic and ran the look-elsewhere search
is available as supplementary material~\cite{phase2_audit}. Material
cut from earlier versions of this paper (interpretive readings of the
ingredients, the circulant analysis, a spectral-determinant near-miss,
and the residual after the cubic) is preserved in the project
repository at \texttt{docs/chronicles/paper2\_narrowing\_cut.md}.
Computational implementation and documentation drafting were performed
with the assistance of large language models (Anthropic Claude) under
the author's scientific direction.
\end{acknowledgments}

\begin{thebibliography}{20}

\bibitem{feynman1985}
R.~P. Feynman, \textit{QED: The Strange Theory of Light and Matter}
(Princeton University Press, 1985).

\bibitem{eddington1935}
A.~S. Eddington, \textit{New Pathways in Science}
(Cambridge University Press, 1935).

\bibitem{barrow2002}
J.~D. Barrow, \textit{The Constants of Nature}
(Vintage Books, 2002).

\bibitem{paper7}
J.~Loutey, ``The Dimensionless Vacuum: Recovering the Schr\"odinger
Equation on a Scale-Invariant Label Lattice,'' GeoVac Paper~7 (2026).

\bibitem{hopf1931}
H.~Hopf, ``\"Uber die Abbildungen der dreidimensionalen Sph\"are auf
die Kugelfl\"ache,'' \textit{Math.\ Ann.}\ \textbf{104}, 637 (1931).

\bibitem{fock1935}
V.~Fock, ``Zur Theorie des Wasserstoffatoms,''
\textit{Z.\ Phys.}\ \textbf{98}, 145 (1935).

\bibitem{alpha_hopf_v1}
% v1 title, cited as superseded [retracted 2026-11-07: g5-retired-titles-2026-11-07]
J.~Loutey, ``The Fine Structure Constant as a Spectral Coincidence
on the Hopf Fibration,'' GeoVac Paper~2, v1 (2026). [Superseded by
the present work.]

\bibitem{wyler1969}
A.~Wyler, ``L'espace sym\'etrique du groupe des \'equations de
Maxwell,'' \textit{C.\ R.\ Acad.\ Sci.\ Paris}\ \textbf{269A}, 743
(1969).

\bibitem{robertson1971}
B.~Robertson, ``Wyler's expression for the fine-structure constant,''
\textit{Phys.\ Rev.\ Lett.}\ \textbf{27}, 1545 (1971).

\bibitem{phase2_audit}
J.~Loutey, ``Paper~2 Computational Audit: Phases 1--4,'' GeoVac
supplementary material (2026). The combinatorial-search driver behind the $p$-value is preserved
durably at \texttt{tests/paper2\_support/phase2\_combinatorial.py}
(2026-07-04; archived-measured --- the full $1.92\times10^9$ sweep is
not regression-gated).

\bibitem{phase6_hopf}
J.~Loutey, ``Quaternionic Hopf Fibration \& Universal Algebraic
Identity: Phase~6,'' GeoVac supplementary material (2026).

\bibitem{alpha_program_record}
J.~Loutey, record of the structural-decomposition program for
Paper~2, Phases 4B--4I (2026), in the GeoVac repository history at
version v2.10.0 (commit \texttt{be9d59d}); Sprint~K-CC in
\texttt{CHANGELOG.md}, v2.27.2.

\bibitem{alpha_sprint_a_memo}
J.~Loutey, ``Paper~2 Sprint~A: Structural consistency verification,''
Sprint record in \texttt{CHANGELOG.md} (2026).

\end{thebibliography}

\end{document}
